\begin{mistake}{2026-04-09}{09}

\begin{problemcontent}
求数列极限：\(\displaystyle I = \lim_{n \to \infty} n \tan(\pi \sqrt{n^2+1}) \)
\end{problemcontent}

\begin{solutioncontent}
利用正切函数的周期性 \(\tan(x) = \tan(x - n\pi)\) 将变量转换为无穷小量：
\[
\begin{aligned}
\tan(\pi \sqrt{n^2+1}) 
&= \tan(\pi \sqrt{n^2+1} - n\pi) \\
&= \tan\left( \pi (\sqrt{n^2+1} - n) \right)
\end{aligned}
\]
对括号内进行分子有理化：
\[ \sqrt{n^2+1} - n = \frac{(\sqrt{n^2+1} - n)(\sqrt{n^2+1} + n)}{\sqrt{n^2+1} + n} = \frac{1}{\sqrt{n^2+1} + n} \]
当 \(n \to \infty\) 时，\(\frac{\pi}{\sqrt{n^2+1} + n} \to 0\)。
利用等价无穷小 \(\tan x \sim x\) \((x \to 0)\)：
\[
\begin{aligned}
I &= \lim_{n \to \infty} n \cdot \tan\left( \frac{\pi}{\sqrt{n^2+1} + n} \right) \\
&= \lim_{n \to \infty} n \cdot \frac{\pi}{\sqrt{n^2+1} + n} \\
&= \lim_{n \to \infty} \frac{\pi}{\sqrt{1+\frac{1}{n^2}} + 1} \\
&= \frac{\pi}{2}
\end{aligned}
\]
\end{solutioncontent}

\mistakeknowledge{
\begin{itemize}
 \item \textbf{核心公式/定理}：等价无穷小替换 \(\tan x \sim x\)；分子有理化技巧。
 \item \textbf{易错点}：直接使用等价无穷小导致逻辑错误（\(\infty\) 时不能用）；忽略三角函数周期性去凑无穷小。
 \item \textbf{关联知识}：凡遇 \(\lim f(n)\) 且含有三角函数时，若自变量趋于无穷，首选利用周期性将其平移至趋于 \(0\) 处。
\end{itemize}}

\end{mistake}
